\documentclass[11pt,a4paper]{amsart}
\usepackage[hmargin=.83in,vmargin=.72in]{geometry}
\usepackage[T1]{fontenc}
\usepackage{amssymb}
\usepackage[hidelinks]{hyperref}
\hypersetup{pdftitle={A sharp beta-density inequality for PAC mode estimation}}
\urlstyle{same}
\setlength{\emergencystretch}{2em}
\linespread{0.95}
\DeclareMathSizes{8}{8}{8}{8}
\DeclareMathSizes{9}{9}{8}{8}
\DeclareMathSizes{10.95}{10.95}{8}{8}
\renewcommand{\bibliofont}{\normalsize\raggedright}
\newtheorem{theorem}{Theorem}
\newtheorem{corollary}[theorem]{Corollary}
\title[A sharp beta-density inequality]{A sharp beta-density inequality for PAC mode estimation}
\date{}
\begin{document}
\begin{abstract}
Anand Jain et al.\ conjectured a lower bound for the common height where two beta densities cross. The bound would imply that the PPR-1v1 stopping rule never stops later than PPR-1vr on the same sample path. We prove the density comparison with sharp factor one, replacing the conjectured factor $(K-1)/K$, and include a zero runner-up count. The proof moves a common trial parameter while holding both counts fixed. A logarithmic derivative gives a quantitative improvement when there are other observations.
\end{abstract}
\maketitle

\section{The crossing and its binary baseline}
For integers $t\ge0$ and $0\le k\le t$, let
\[
 L_{t,k}(p)=\frac{p^k(1-p)^{t-k}}{\mathrm{B}(k+1,t-k+1)}
          =(t+1)\binom{t}{k}p^k(1-p)^{t-k}.
\]
At endpoints the formula is interpreted by continuity. Fix integers $a>b\ge0$ and $f\ge0$, and put $c=a+b$ and $t=c+f$. The leader and competitor have counts $a,b$; $f$ is the remaining count. Their unique interior crossing is
\[
 r=\left(\frac{a!(t-a)!}{b!(t-b)!}\right)^{1/(a-b)},
 \qquad \theta=\frac{r}{1+r}.
\]
The formula follows from the density ratio on $(0,1)$. The common height will be written $H=L_{t,a}(\theta)=L_{t,b}(\theta)$. At $f=0$ the crossing is $1/2$, with common height $H_0=L_{c,a}(1/2)=L_{c,b}(1/2)$.

Appendix~C.2 of Anand Jain et al.\ \cite[PDF p.~19]{aj} conjectures $H\ge (K-1)H_0/K$ for the mode-estimation comparison with $K$ categories. The proposed bound was checked numerically on positive counts. The following bound supplies the required implication with a sharp constant and covers $b=0$.

\begin{theorem}\label{thm:bound}
For all the above counts,
\[
     H\ \ge\ \frac{(c+f+1)^2}{(c+1)(c+2f+1)}\,H_0\ \ge\ H_0,
 \qquad H_0=\frac{1}{2^c\,\mathrm{B}(a+1,b+1)}.
\]
Equality in $H\ge H_0$ holds exactly when $f=0$.
\end{theorem}

\section{A moving trial point}
Put $x=t+1$ and $m=(c+1)/2$. The trial point is $p=m/x$. For an integer $k\in[0,c]$ define a positive real interpolation for all $x\ge c+1$ by
\[
 F_k(x)=\frac{\prod_{j=0}^k(x-j)}{k!}
          \left(\frac{m}{x}\right)^k
          \left(1-\frac{m}{x}\right)^{x-1-k}.
\]
For integer $x$ it equals $L_{x-1,k}(m/x)$. A logarithmic derivative gives
\[
 \frac{\mathrm d}{\mathrm dx}\log F_k(x)
 =\sum_{j=0}^{k}\frac1{x-j}-\frac{k}{x}
   +\log(1-m/x)
   +\frac{(x-1-k)m}{x(x-m)}.
\]
The summands for $j=1,\ldots,k$ are right endpoint values of the increasing function $s\mapsto 1/(x-s)$. Hence the sum is at least $1/x-\log(1-k/x)$, including $k=0$. Write $\alpha=k/x$, $\beta=m/x$ and $z=(1-\alpha)/(1-\beta)$. Regrouping the derivative bound yields
\[
 \frac{\mathrm d}{\mathrm dx}\log F_k(x)
 \ \ge\ z-1-\log z+\frac{1-2\beta}{x(1-\beta)}
 \ \ge\ \frac{x-2m}{x(x-m)}.
\]
Here $\log z\le z-1$ for $z>0$. Integrating the last expression from $2m$ to $x$ gives
\[
 F_k(x)\ \ge\ F_k(2m)\frac{x^2}{4m(x-m)}.
\]
For $k=a,b$, both initial values equal $H_0$ by binomial symmetry. The factor is one at $x=2m$ and larger for $x>2m$, since $x^2-4m(x-m)=(x-2m)^2$.

\section{From the trial to the crossing}
The crossing odds are the geometric mean
\[
 r=\left(\prod_{j=b+1}^{a}\frac{j}{t-j+1}\right)^{1/(a-b)}.
\]
The factors increase in $j$. Bounding their geometric mean by the first and last factor gives $b/t\le \theta\le a/t$. For $p\in(0,1)$ the ratio $L_{t,a}(p)/L_{t,b}(p)$ increases strictly and equals one at $\theta$. On the left of $\theta$, the smaller density is $L_{t,a}$, which increases up to its mode $a/t$. On the right it is $L_{t,b}$, which decreases from its mode $b/t$. These monotonicities extend to boundary modes $0$ and $1$ by continuity. Therefore the crossing maximizes the smaller density on all of $[0,1]$. In particular
\[
       H\ \ge\ \min\bigl\{F_a(t+1),F_b(t+1)\bigr\}.
\]
The integrated bound gives Theorem~\ref{thm:bound}. At $f=0$, $r=1$ and both densities meet the binary baseline, so the universal factor one is sharp.

For instance, with $a=2$, $b=0$, $f=7$, the crossing is $1/7$ and its height is $10(6/7)^9$. The common trial is $3/20$. With the same $a,b$ and $f=0$, the two densities meet at $1/2$ with height $3/4$. The illustration shows the lower envelope in these examples.

\section{Pathwise consequence for the original PPR rules}
\begin{corollary}
For the original uniform-prior PPR-1v1 and PPR-1vr rules of \cite{aj}, with the same error parameter $0<\delta<1$ and $K\ge2$ categories, PPR-1v1 stops no later than PPR-1vr on each fixed sample path.
\end{corollary}
\begin{proof}
In the 1vr rule the full-count interval for a category with $k$ observations is the superlevel interval of $L_{t,k}$ at threshold $\delta/K$ \cite[Sections~3.2 and C.2]{aj}. Its interior in $(0,1)$ includes every point of strict superlevel density there. Tied leaders cannot be separated, because their intervals coincide and contain a positive-length strict superlevel interval. The latter follows from normalization and $\delta/K<1$.

If the 1vr rule stops, take its empirical leader of count $a$ and any competitor of count $b<a$. Their interval interiors are disjoint, so both densities cannot strictly exceed $\delta/K$ at their crossing. Hence $H\le\delta/K$. The corresponding pair-only 1v1 test uses $H_0$ at $1/2$ and threshold $\delta/(K-1)$. The theorem yields
\[
                H_0\le H\le \delta/K < \delta/(K-1).
\]
Thus the leader passes all pairwise tests by this time, including against zero-count competitors. The strict comparison also covers 1vr intervals touching at a threshold endpoint.
\end{proof}

Qiu, Artola Velasco and Gomez-Rodriguez \cite[Appendix~A.1]{qiu} vary the counts while evaluating at $1/2$. Here the total varies with both counts fixed.

\section*{Replay and technical disclosure}
The standard-library replay uses exact rational arithmetic for the trial bound and crossing-odds brackets on a finite parameter family. It also checks the displayed rational crossing fixture. External peer review and formal verification are pending.

ChatGPT Astra 6.2 (OpenAI) generated the recovered proof and original replay under human direction. OpenAI Codex independently reviewed the mathematics. GPT-6.2 Sol (OpenAI) checked primary sources and prepared the manuscript, extended replay and illustration. DannyExperiments is the project citation identity.

\bibliographystyle{amsplain}
\bibliography{references}
\end{document}
